\section{Proof of the main theorem}\label{app:quantitative}
We deduce both bounds of Theorem~\ref{thm:main} from Boolean tensor
contraction. A formula of fixed clause density has a network with
$O_{k,D}(n)$ wires. For arbitrary density, disjoint sparsification and
parity tests reduce the count to sparse polynomial systems, whose
networks have $O_k(n\log\log n)$ wires. We state the needed outputs
first and give the two deductions in
Sections~\ref{subsec:sharp-fixed-density} and~\ref{subsec:sharp-dense}.

\subsection{Boolean tensor contraction}
\label{subsec:tensor-interface}

A Boolean tensor network is a bipartite multigraph $G$ in which each
edge, also called a wire, carries a single bit. Each vertex $v$ has an ordered list
$\delta(v)$ of incident edges and a tensor $T_v$ assigning a ring
element to each Boolean assignment to those edges. Parallel edges
have distinct edge ends. The \emph{contraction} of the network is
\[
 Z=\sum_{x\in\{0,1\}^{E(G)}}\prod_{v\in V(G)} T_v(x|_{\delta(v)}).
\]
An isolated vertex contributes its scalar entry. We measure the size
of the network by its number of edges.

\begin{theorem}[Boolean tensor contraction]\label{thm:tensor-contraction}
Fix positive constants $C,A,b$.  Let $s\ge2$ be an integer, and let
$\mathcal L$ be the total length of the graph and tensor descriptions
of a bipartite tensor network with at most $Cs$ edges and maximum
degree at most $\lceil s^{9/10}\rceil$. Let $P_0,\ell$ be positive
integers. Suppose that the number of
vertices is bounded by a fixed polynomial in $s$, and that the entries
belong to
\[
 \mathcal B_{P_0,\ell}
   =(\mathbb Z/2^{P_0}\mathbb Z)[V]/(V^\ell),
 \qquad 1\le P_0,\ell\le Cs.
\]
Suppose, uniformly for every vertex and Boolean choice of its edge ends,
that its entry can be computed from the description in time
$(s+\mathcal L)^b2^{A s^{9/10}}$ and space $(s+\mathcal L)^b$.
Then the contraction can be computed deterministically in
\[
 (s+\mathcal L)^{O(1)}
  2^{O(s/\sqrt{\log s})}
\]
bit operations and $(s+\mathcal L)^{O(1)}$ bits of space.  Constants
depend only on the fixed bounds in the hypotheses.
\end{theorem}

The proof is in Section~\ref{app:tensor}. At recursion depth $d$,
a call that is not a leaf has $M=O(s8^{-d})$ wires and
$u=\Theta((\log s)/2^d)$ retained two-adic valuation levels.
Interactions between groups are expanded in auxiliary variables representing linear forms in column bits.
Only $2^{O(M/\sqrt u)}$ monomials in these auxiliary variables survive the working precision,
and each factors into contractions with at most $M/8$ wires.
The exponents therefore sum geometrically to $O(s/\sqrt{\log s})$.


\subsection{Counting sparse polynomial systems}
The following consequence of contraction provides the counter needed
for every candidate system in the reduction, including candidates
that still accept nonsatisfying assignments.

\begin{lemma}[Counting sparse polynomial systems by contraction]
\label{lem:tensor-sparse-query}
Fix integers $k,c\ge1$. Given polynomials $P_1,\ldots,P_a$ over
$\mathbb F_2$ on $n\ge20$ declared Boolean variables, with
$a\le cn$ and each $P_i$ listed as at most
$\lceil\log_2 n\rceil^c$ monomials of degree at most $k$, the number
of common zeros in $\{0,1\}^n$ is computable deterministically in
$n^{O_{k,c}(1)}2^{O_{k,c}(n\log\log n/\sqrt{\log n})}$ bit operations and
$n^{O_{k,c}(1)}$ bits of space.
\end{lemma}
The proof is in Section~\ref{sec:tensor-encoding}.

\subsection{Sparsification and parity tests}\label{subsec:main-reductions}
Sparsification supplies formulas with few occurrences per variable,
and the parity tests compress each leaf into $O(n)$ equations.

\begin{lemma}\label{lem:sparsify}
Fix $k\ge2$. Let $F$ be a $k$-CNF on $n$ declared variables with
$m$ clauses, and let $0<\sigma\le1$ be rational. Write $L_\sigma$
for the total binary encoding length of its supplied numerator and
denominator. A deterministic algorithm constructs a tree of $k$-CNFs
on the same variables with the following properties.
An assignment satisfies the input exactly when it satisfies one
leaf formula, and it cannot satisfy two distinct leaves.
The tree has
\[
 \operatorname{poly}_k(n,\sigma^{-1})\,2^{\sigma n}
\]
nodes and depth $\poly_k(n,\sigma^{-1})$.
In every leaf, each variable occurs at most $B$ times, where $B$
is an explicitly computed integer satisfying
\[
 B=O_k\!\Bigl(\bigl[\sigma^{-1}\log(2/\sigma)\bigr]^{k-1}\Bigr).
\]
Input normalization uses $\poly_k(n,m)$ bit operations.
Constructing the numerical parameters uses $\poly_k(L_\sigma)$
bit operations. Thereafter, processing each node uses
$\poly_k(n,\log(2/\sigma))$ bit operations.
\end{lemma}
Thus the input count is the sum of the leaf counts.
Appendix~\ref{app:sparsify} proves the lemma, including the disjoint
branching rule and the bounds for a varying rational parameter.

For a system $\mathcal P$ of equations $P_i(x)=0$ on $n$ Boolean
variables, write
\begin{equation}\label{eq:common-zero-count}
 N(\mathcal P)=\#\{x\in\{0,1\}^n:P_i(x)=0\text{ for every }i\}.
\end{equation}
Let $F$ be a $k$-CNF on $n$ declared variables with $m\ge1$
indexed clauses, and let $v_e(x)$ be the violation bit
of clause $e$. As a polynomial over $\mathbb F_2$, $v_e$ is a product
of the factors $1+x_v$ for positive literals and $x_v$ for negative
literals. A tautological clause has $v_e=0$ and an empty clause has
$v_e=1$. After removing repeated literals, $\deg v_e\le k$ and $v_e$
expands to at most $2^k$ monomials. Put $v(x)=(v_1(x),\ldots,v_m(x))$.
For a binary matrix $A$ with $m$ columns, row $i$ specifies the equation
\begin{equation}\label{eq:hash-query}
 P_i(x):=\sum_{e:A_{ie}=1} v_e(x)=0\qquad\text{over }\mathbb F_2.
\end{equation}
Write $\mathcal P_A$ for this system and $N(A)=N(\mathcal P_A)$ for its
number of common zeros. Every satisfying assignment passes every row.
Hence $N(A)=\#F+D(A)$, where $D(A)$ counts the nonsatisfying assignments
still passing. We do not know $\#F$, but it is the same for every
candidate next row: the row minimizing $N$ also minimizes $D$.

We represent a test row by a list of pairs $(e,c)$ with
$e\in[m]$ and $c\in\{0,1\}$, summing coefficients modulo $2$ if a
clause index is repeated. The length of this list counts all pairs,
including repeated indices and pairs with coefficient zero. A row of
list length $s$ thus reads at most $ks$ variables before cancellation.


The initial tests exclude every nonzero violation vector with small
support. They will leave only assignments that a short additional
test catches with constant probability.

\begin{lemma}\label{lem:sparse-checks}
Let $n,m,B$ be positive integers with $m\le Bn$, and put
\[
 \mu=\lfloor\log_2 m\rfloor+1,\qquad
 P=20000B^2\mu^2,\qquad D_0=320000B^2\mu^3.
\]
There is a deterministic construction of a binary matrix
$A^{(0)}$ with $m$ columns and at most $3n$ rows such that every row
has at most $D_0$ nonzero entries and
\[
 z\in\{0,1\}^m\setminus\{\vec{0}\},\quad
 A^{(0)}z=\vec{0}\pmod2
 \quad\Longrightarrow\quad
 |\operatorname{supp}z|\ge m/P.
\]
The construction uses $\exp(\poly(B,\log(m+2)))$
bit operations and $\poly(n,m,B)$ bits of working space.
\end{lemma}

\begin{lemma}\label{lem:greedy-hashes}
Let $F$ be a $k$-CNF on $n$ declared variables with $1\le m\le Bn$
indexed clauses, and take $P,D_0,A^{(0)}$ from
Lemma~\ref{lem:sparse-checks}. Using common-zero counts of the
systems \eqref{eq:hash-query}, a deterministic procedure extends
$A^{(0)}$ by $3n$ rows of list length $P$ to a matrix $A$ satisfying
$N(A)=\#F$. It uses at most $3n(2m)^P$
common-zero counts. Every candidate system has at most $6n$ equations,
each involving at most $2^k\max(D_0,P)$ monomials of degree at most $k$.
\end{lemma}
Section~\ref{sec:hashes} proves these two lemmas.

\subsection{Fixed clause density}\label{subsec:sharp-fixed-density}

\begin{proof}[Proof of the fixed-density bound in Theorem~\ref{thm:main}]
Give each clause a vertex carrying its Boolean satisfaction table, with
one edge-end for each distinct variable in the clause. A tautology may be
dropped, and an empty clause gives count zero. For a variable with
$d\ge1$ edge ends, require their bits to agree. For $d=1$ attach the unary
tensor $(1,1)$, and for $d=2$ attach a binary equality tensor. For
$d\ge3$ take a tree with $d$ external edge ends and trivalent equality
tensors at its internal vertices. Such a tree has $d-2$ internal
vertices, and each common value on its edge ends extends uniquely to all
its internal edges.

There are at most $km$ clause edge ends. The graph has $O_k(n+m)$ edges
and vertices, and maximum degree at most $\max\{k,3\}$. Subdivide
every edge once by a binary equality tensor. The old vertices and the
new vertices form the two sides of a bipartite network. The two bits
on a subdivided edge must be equal, so this preserves the contraction.
Every satisfying assignment on the occurring variables contributes
one and every other assignment contributes zero.

With $s=n+m$, apply Theorem~\ref{thm:tensor-contraction} over
$\mathbb Z/2^{n+1}\mathbb Z$, taking $P_0=n+1$ and $\ell=1$. The
local tables have bounded size for fixed $k$. The
degree and precision bounds hold for all sufficiently large $s$, and
bounded instances can be enumerated directly. If $u$ declared
variables are unused, multiply the contraction by $2^u$ modulo
$2^{n+1}$. Its least nonnegative residue is $\#F$, since
$0\le\#F\le2^n$. This includes the empty formula, for which the
empty network has contraction one. The time is
$s^{O_k(1)}2^{O_k(s/\sqrt{\log s})}$ and the space is polynomial.
When $m\le Dn$, we have $n\le s\le(1+D)n$, giving exponent
$O_{k,D}(n/\sqrt{\log n})$.
\end{proof}


\subsection{Arbitrary clause density}\label{subsec:sharp-dense}
\begin{proof}[Proof of the arbitrary-density bound in
Theorem~\ref{thm:main}]
Normalize the formula by deleting repeated literals, tautological
clauses, duplicate clauses, and clauses containing another clause.
Retain all $n$ declared variables. Empty clauses and
formulas with no clauses have immediate answers. For $n\ge20$
apply Lemma~\ref{lem:sparsify} with
$\sigma=1/\lceil\sqrt{\lambda}\rceil$, where
$\lambda=\lceil\log_2 n\rceil$. The tree has
$n^{O_k(1)}2^{O(n/\sqrt{\lambda})}$ nodes and polynomial depth.
Every leaf has occurrence bound
$B=O_k((\sigma^{-1}\log(2/\sigma))^{k-1})=\lambda^{O_k(1)}$
and at most $Bn$ nonempty clauses.

For a nontrivial leaf with $m_G$ clauses, use the initial parity
checks of Lemma~\ref{lem:sparse-checks} and the greedy procedure
of Lemma~\ref{lem:greedy-hashes}. Each initial row has support size
at most $D_0$, and each additional row has list length $P$.
Both bounds are $\lambda^{O_k(1)}$. By \eqref{eq:hash-query}, every row
gives a degree-at-most-$k$ equation with at most
$2^k\max\{P,D_0\}$ monomials, and every queried system has at
most $6n$ equations. A single constant $c=c(k)$ bounds both the
ratio of equations to variables and the exponent of $\lambda$ in
the monomial bound for all sufficiently large $n$.
At each greedy step, use Lemma~\ref{lem:tensor-sparse-query} to
compute $N(A;h_\gamma)$ for every candidate row $h_\gamma$, where
$h_\gamma$ is the row represented by the list $\gamma$ and
$A;h_\gamma$ is $A$ with that row appended. Append a candidate
attaining the minimum. The evaluator applies also to
discarded candidates whose systems still accept nonsatisfying
assignments. One final call gives $N(A)$ for the selected matrix,
which equals the leaf count by Lemma~\ref{lem:greedy-hashes}.

Constructing the initial checks and enumerating candidates costs
$2^{\lambda^{O_k(1)}}$ bit operations apart from the evaluations.
There are at most $3n(2m_G)^P+1=2^{\lambda^{O_k(1)}}$ evaluations
per leaf. Since every fixed power of $\lambda$ is
$o(n\log\log n/\sqrt{\log n})$, these costs and the number of leaves
are absorbed in the claimed time bound
\[
 (n+m)^{O_k(1)}2^{O_k(n\log\log n/\sqrt{\log n})}.
\]
Here $m$ counts the input clauses. Traverse the tree
depth first and process one query at a time. The leaf solution sets
are disjoint, so their exact counts add with an $(n+1)$-bit
accumulator. Together with polynomial space per query and per
tree path, this gives the asserted polynomial-space bound.
\end{proof}

A single algorithm attains both bounds. With
$\lambda=\lceil\log_2 n\rceil$ and $\nu=\lceil\log_2\lambda\rceil$,
use the clause network when $m\le n\nu$ and the parity queries
otherwise. In the first branch $n\le s=n+m=O(n\log\log n)$,
so its exponent is $O_k(n\log\log n/\sqrt{\log n})$ even at growing
density. For each fixed $D$ and all sufficiently large $n$, every
input with $m\le Dn$ takes that branch with $s=O_D(n)$.
